\begin{titlepage}

    \begin{center}
        \includegraphics[height=0.45\textheight]{2026-03-06_Bat_Source_Multiboard_Render.png}
    \end{center}
    \todo[inline]{Render durch Foto des fertigen Geräts im Gehäuse ersetzen}

    \vspace{2mm}
    \begin{center}
        \rule{\textwidth}{2pt}
    \end{center}
    \vspace{2mm}

    \begin{center}
        \huge{\textbf{\ST}}\\
        \vspace{4mm}
        \LARGE{\textbf{Batteriebetriebenes Prüfgerät für die elektrische Abnahme in der Formula Student}}\\
        \vspace{4mm}
        \Large{\ManualName}
    \end{center}

    \vspace{2mm}
    \begin{center}
        \rule{\textwidth}{2pt}
    \end{center}
    \vspace{4mm}

    \begin{center}
        \large{\makeatletter \@author \makeatother}\\
        \vspace{6mm}
        Revision Anleitung: \tbc{0.1} \quad Firmware: \tbc{x.y} \quad Hardware: \tbc{ST1-1BD}\\
        \vspace{2mm}
        \today
    \end{center}

    \vfill
    \begin{warning}
        Diese Anleitung vor dem Gebrauch vollständig lesen, insbesondere \autoref{sec:safety}.
        Das Gerät kann bis zu \qty{1000}{\volt} erzeugen.
    \end{warning}

\end{titlepage}
